%% parts/p04/frontiers-of-computability.tex
%% Frontiers survey with cited results and proved Omega/computable-structure results.

\chapter{Frontiers of Computability}
\label{sec:p04-frontiers-computability}

\noindent\textit{Reader note.} This is a separately authored survey following
the exact inherited computability chapters.  Its hierarchy, information,
Diophantine, word-problem, tiling, Busy-Beaver, and computable-structure
statements are controlled by the displayed sources and do not replace any
inherited theorem or proof.

The equivalence of classical formal models and the Church--Turing thesis have
different standings.  The former compares exact classes of numerical
functions; the latter identifies one such class with the informal notion of
effective calculability.  The distinction fixed in
Boundary~\ref{def:p04-church-turing} remains in force here
\cite{Turing1937Lambda,Soare1996ComputabilityRecursion}.  This chapter studies
mathematical boundaries obtained after a formal model, an input encoding, and
a yes predicate have been fixed.

\section{The lightface analytical and hyperarithmetical hierarchies}
\label{sec:p04-analytical-hierarchy}

\begin{hbconvention}{def:p04-analytical-language}{Second-order arithmetic convention}
Number variables range over \(\N\), and function variables range over
\(\N^{\N}\).  Equivalently, a function may be replaced by the subset of
\(\N\times\N\) coding its graph, and then by a subset of \(\N\) through the
fixed numerical pairing.  A matrix is \emph{arithmetical} when, after its displayed
function variables are treated as function or predicate symbols, it has only
number quantifiers.  Throughout this section the word \emph{analytical} means
\emph{lightface}: no function or set parameter is allowed.  Finite numerical
parameters may be absorbed into the effective code.
\index{analytical hierarchy!lightface}
\index{arithmetical matrix}
\end{hbconvention}

\begin{hbdefinition}{def:p04-lightface-analytical-hierarchy}{Lightface analytical hierarchy}
For \(n\geq1\), a relation \(A\subseteq\N^k\) belongs to
\(\Sigma^1_n\) when there is an arithmetical formula \(\varphi\) such that
\[
 \vec x\in A\quad\Leftrightarrow\quad
 \begin{cases}
  \exists\alpha_1\,\forall\alpha_2\,\exists\alpha_3\cdots
    \exists\alpha_n\,\varphi(\vec x,\alpha_1,\ldots,\alpha_n),
       & n\text{ odd},\\
  \exists\alpha_1\,\forall\alpha_2\,\exists\alpha_3\cdots
    \forall\alpha_n\,\varphi(\vec x,\alpha_1,\ldots,\alpha_n),
       & n\text{ even}.
 \end{cases}
\]
The displayed function quantifiers form \(n\) alternating blocks and the first
block is existential; the parity clauses make the final quantifier explicit.
The class \(\Pi^1_n\) is defined by a prefix
beginning universally, and
\[
 \Delta^1_n\defeq\Sigma^1_n\cap\Pi^1_n.
\]
At level zero, \(\Sigma^1_0=\Pi^1_0=\Delta^1_0\) is the class of
arithmetical relations.  These symbols name effective pointclasses of
relations, not syntactic formula classes and not the boldface projective
pointclasses \cite{Kleene1955Hierarchy,Moschovakis2016Hyperarithmetical}.
\index{Sigma-1-n@\(\Sigma^1_n\)}
\index{Pi-1-n@\(\Pi^1_n\)}
\index{Delta-1-n@\(\Delta^1_n\)}
\end{hbdefinition}

\begin{hbdefinition}{def:p04-universal-analytical-sets}{Universal analytical sets}
For each \(n\geq1\), effectively enumerate the valid lightface
\(\Sigma^1_n\) prefix-form definitions in the fixed second-order arithmetic
syntax.  Put
\[
 \mathsf U^1_{\Sigma_n}
 =\{\langle e,x\rangle:\text{definition }e\text{ holds of }x\},
 \qquad
 \mathsf U^1_{\Pi_n}=\N^2\setminus\mathsf U^1_{\Sigma_n},
\]
with invalid definition codes interpreted as false on the \(\Sigma\)-side.
These are universal coded relations, not additional pointclasses.  Their index
coordinate records a definition, unlike a fixed natural parameter inside one
lightface formula.  The distinction is needed to provide one calibrated
complete relation at every analytical level in the table below.
\index{universal analytical set}
\end{hbdefinition}

\begin{hbtheorem}{thm:p04-universal-analytical-completeness}{cited=Kleene1955Hierarchy}{Complete relations at every analytical level}
For every \(n\geq1\), \(\mathsf U^1_{\Sigma_n}\) is
\(\Sigma^1_n\)-complete and \(\mathsf U^1_{\Pi_n}\) is
\(\Pi^1_n\)-complete under effective many-one reductions.
\end{hbtheorem}

\begin{hbmethod}{def:p04-kleene-ordinal-notations}{Kleene ordinal notations and iterated jumps}
Fix Kleene's standard cited system \((\mathcal O,<_{\mathcal O})\), rather
than treating the following orientation clauses as an independent eliminable
definition.  The base code \(1\) denotes \(0\).  If \(a\in\mathcal O\), then
\(2^a\in\mathcal O\) is a successor notation.  If \(e\) indexes a total
computable sequence \(a_0<_{\mathcal O}a_1<_{\mathcal O}\cdots\) of notations,
then \(3\cdot5^e\in\mathcal O\) is a notation for their supremum.
Here \(<_{\mathcal O}\) compares notation codes in the fixed standard system;
it is not simply the order between their denoted ordinals, since distinct
notations can denote the same recursive ordinal.

The associated effective jump hierarchy starts with \(H_1=\varnothing\),
puts \(H_{2^a}=H_a'\), and at an effective limit takes the computable join of
the earlier sets named by the chosen sequence.  In this chapter
\[
 \mathsf{HYP}\defeq
 \{A\subseteq\N:\text{for some }a\in\mathcal O, A\leq_T H_a\}.
\]
This is a coding-dependent notation system used to organize effective
transfinite jump iteration, not a canonical injective enumeration of all
countable ordinals. The simultaneous construction, hierarchy existence and
characterization are supplied in
Section~\ref{sec:p04-delta-one-hyperarithmetical}. The cited historical
presentations remain separately identified \cite{Kleene1955Hierarchy,
Moschovakis2016Hyperarithmetical,Rogers1967}.
\index{Kleene O@Kleene's \(\mathcal O\)}
\index{hyperarithmetical set}
\end{hbmethod}

\begin{hbtheorem}{thm:p04-hyp-delta-one-one}{proved=annex:p04.doh.overview}{Hyperarithmetical characterization}
For subsets of \(\N\), under the effective codings fixed above,
\[
 \mathsf{HYP}=\Delta^1_1.
\]
Equivalently, a set is hyperarithmetical exactly when it is both
\(\Sigma^1_1\) and \(\Pi^1_1\).  This is a lightface theorem; it does not
identify \(\mathsf{HYP}\) with either one of those two larger pointclasses.
\end{hbtheorem}

% Independent O/jump construction; not a replacement for a historical source.
\subsection{Effective transfinite computation and delta-one-one definitions}
\label{sec:p04-delta-one-hyperarithmetical}

There are two descriptions of hyperarithmetical sets. One iterates ordinary
oracle jumps along effective ordinal notations. The other requires a set and
its complement to have existential second-order definitions with arithmetic
matrices. The equivalence is a construction theorem: a complementary pair of
definitions can be compiled into a single valid jump oracle and a total
decision program. The following proof supplies that compilation for the
\(\mathcal O\) convention of
Definition~\ref{def:p04-kleene-ordinal-notations}.
It uses neither an assumed Church--Kleene rank bound nor the equality it
proves. Effective grounded recursion and the distinction between ordinal
systems are discussed in \cite[\S\S2A--3E]{Moschovakis2016HYPAuthorManuscript};
the exact \(\mathcal O\) clauses are also displayed in
\cite[\S2.3]{Speranski2017TruthAuthorPreprint}.

\subsubsection{Programs, notations and actual jump sets}

\begin{hblemma}{thm:p04-doh-arithmetic-graphs}{proved=annex:p04.doh.arithmetic-graphs}{Arithmetic formulas for computable numerical graphs}
Every partial computable numerical function has its graph defined by an
arithmetical formula in the standard model with vocabulary
\(0,1,+,\times\), obtained effectively from a machine for the function.
Every computable predicate, finite-word operation and substitution function
therefore has an arithmetic definition. For a total function its graph
has exactly one output at each standard-model input.
\end{hblemma}

\begin{hbconvention}{def:p04-doh-programs}{Arithmetic and oracle coding}
Work over the standard natural numbers in classical ZF set theory. Fix a
computable pairing bijection \(\langle i,x\rangle\) and effective finite
sequence coding. Fix an acceptable numbering \(\varphi_e^B\) of partial
oracle programs: universal simulation and insertion of numerical parameters
are computable operations on program codes, independent of the oracle.
For a set \(B\subseteq\N\), put
\[
 B'=\{e:\varphi_e^B(e)\text{ is defined}\}.
\]
Definedness uses the numerical-output convention: rejection or a physical
halt without a valid numerical output does not define a value. This is
the same Turing-jump convention as in the earlier oracle presentation.
A reduction \(A\leq_T B\) computes the total characteristic function of
\(A\), with values \(0,1\), from oracle \(B\).
For sets \(B_i\), their effective join is
\[
 \bigoplus_i B_i=\{\langle i,x\rangle:x\in B_i\}.
\]
An index compiler constructs a finite program; it need not run the program
or decide whether a prospective ordinal notation is valid. Computable
numerical operations are expressed by their arithmetic graph definitions;
the independent finite-history construction is supplied in
Theorem~\ref{thm:p04-doh-arithmetic-graphs}. It does not depend on the
hyperarithmetical characterization being proved here.
\end{hbconvention}

\begin{hblemma}{thm:p04-doh-program-tools}{proved=annex:p04.doh.program-tools}{Uniform program self-reference and jump simulation}
For every total computable transformation \(f\) of program codes there is
an index \(e\) satisfying \(\varphi_e^B=\varphi_{f(e)}^B\) for every oracle
\(B\). Numerical parameters can remain arguments of this one program.
There are total computable index transformations witnessing
\(B\leq_T B'\), and taking a reduction \(A\leq_T B\) to a reduction
\(A'\leq_T B'\). Their correctness is conditional only on the supplied
reduction actually being total and correct.
\end{hblemma}

\begin{hbdefinition}{def:p04-doh-notations}{The simultaneous notation system}
Write \(\mathsf s(a)=2^a\), and \(\mathsf l(e)=3\cdot5^e\).
Define \(\mathcal O\) and its predecessor relation \(<_{\mathcal O}\)
by the least simultaneous closure under these clauses:
\begin{enumerate}
 \item \(1\in\mathcal O\).
 \item If \(a\in\mathcal O\), then \(\mathsf s(a)\in\mathcal O\),
 and its predecessors are \(a\) and the predecessors of \(a\).
 \item If \(\varphi_e\) is total, all \(a_i=\varphi_e(i)\) belong to
 \(\mathcal O\), and \(a_i<_{\mathcal O}a_{i+1}\) for every \(i\),
 then \(\mathsf l(e)\in\mathcal O\). Its predecessors are exactly the
 \(x\) such that \(x\leq_{\mathcal O}a_i\) for some \(i\), where
 \(\leq_{\mathcal O}\) includes equality.
\end{enumerate}
``Least'' means that no notation or predecessor pair is added except by
iteration of these clauses. A limit clause is an infinite-premise clause;
it is not a computable validity test. Equivalently, take the least fixed
point of the positive operator on tagged notation codes and predecessor
pairs specified by the clauses, retaining already added records at each
step. Set-sized ordinal iteration constructs this fixed point.

The ordinal denotation is
\[
 \nu_{\mathcal O}(1)=0,\qquad
 \nu_{\mathcal O}(\mathsf s(a))=\nu_{\mathcal O}(a)+1,\qquad
 \nu_{\mathcal O}(\mathsf l(e))=\sup_i\nu_{\mathcal O}(\varphi_e(i)).
\]
The associated sets, with the manuscript's empty-base convention, are
\[
 H_1=\varnothing,\qquad H_{\mathsf s(a)}=H_a',\qquad
 H_{\mathsf l(e)}=\bigoplus_i H_{\varphi_e(i)}.
\]
These are definitions on valid notations. Values on invalid codes are not
silently assigned. This is the smaller \(\mathcal O\) system, not the
system which admits every sequence merely increasing in denoted ordinal.
\end{hbdefinition}

\begin{hblemma}{thm:p04-doh-notation-existence}{proved=annex:p04.doh.notation-existence}{Well-defined denotations and hierarchy}
Every valid notation has a unique ordinal denotation and a unique set
\(H_a\). If \(b<_{\mathcal O}a\), then
\(\nu_{\mathcal O}(b)<\nu_{\mathcal O}(a)\).
In particular, repeatedly following a successor predecessor or a selected
limit-sequence term cannot continue infinitely from a valid notation.
\end{hblemma}

\begin{hbdefinition}{def:p04-doh-addition}{A syntactic notation-addition compiler}
The computable compiler \(\operatorname{add}_{\mathcal O}(a,b)\) obeys
\[
\begin{aligned}
 \operatorname{add}_{\mathcal O}(a,1)&=a,\\
 \operatorname{add}_{\mathcal O}(a,\mathsf s(c))
   &=\mathsf s(\operatorname{add}_{\mathcal O}(a,c)),\\
 \operatorname{add}_{\mathcal O}(a,\mathsf l(e))
   &=\mathsf l(\operatorname{index}(i\mapsto
          \operatorname{add}_{\mathcal O}(a,\varphi_e(i)))).
\end{aligned}
\]
The last line outputs an index without running its sequence. On codes of
none of these syntactic forms put \(\operatorname{add}_{\mathcal O}(a,b)=1\).
Program self-reference fixes the sequence index. Successor parsing recurses
only to the smaller natural number \(c\); hence the top-level compiler is
total even on invalid codes. This totality does not assert validity of its
output on those codes.
\end{hbdefinition}

\begin{hblemma}{thm:p04-doh-addition}{proved=annex:p04.doh.addition}{Notation addition and uniform right-summand access}
For \(a,b\in\mathcal O\), the compiled code belongs to \(\mathcal O\) and
\[
 \nu_{\mathcal O}(\operatorname{add}_{\mathcal O}(a,b))
     =\nu_{\mathcal O}(a)+\nu_{\mathcal O}(b).
\]
It preserves the predecessor relation in its second argument. Also
\(a\leq_{\mathcal O}\operatorname{add}_{\mathcal O}(a,b)\), with a strict
inequality when \(b\ne1\).
There is a total computable index function \(r(a,b)\) such that
\[
 \varphi_{r(a,b)}^{H_{\operatorname{add}_{\mathcal O}(a,b)}}
       =\chi_{H_b}.
\]
The reduction does not infer Turing reducibility from ordinal inequality;
its successor simulation and limit-column program are explicitly constructed.
\end{hblemma}

\subsubsection{An effective oracle bound for a well-founded computation}

\begin{hbdefinition}{def:p04-doh-node-bounds}{Node notations for a recursive child relation}
Let \(\mathsf C(z,i)\) be a decidable relation saying that the natural
number \(i\) is a child of node \(z\). Assume that there is no infinite
sequence \(z_0,z_1,\ldots\) with \(\mathsf C(z_j,z_{j+1})\).
Fix a self-referential computable code generator \(h\) as follows. Its
sequence program at node \(z\) computes
\[
\begin{aligned}
 a_{z,0}&=1,\\
 b_{z,i}&=\begin{cases}h(i),&\mathsf C(z,i),\\1,&\text{otherwise},\end{cases}\\
 a_{z,i+1}&=\operatorname{add}_{\mathcal O}
                  (\mathsf s(a_{z,i}),b_{z,i}),\\
 e_z&=\operatorname{index}(i\mapsto a_{z,i}),\\
 h(z)&=\mathsf s(\mathsf s(\mathsf l(e_z))).
\end{aligned}
\]
To output \(h(z)\), the generator only constructs \(e_z\); it does not
evaluate its children. Each individual sequence value uses finitely many
calls to the total top-level generator and notation-addition compiler.
\end{hbdefinition}

\begin{hblemma}{thm:p04-doh-node-bounds}{proved=annex:p04.doh.node-bounds}{Valid node bounds with uniform child access}
The generator is total computable. Every \(h(z)\) is a valid notation.
Put \(J_z=H_{\mathsf l(e_z)}=\bigoplus_i H_{a_{z,i}}\), so
\(H_{h(z)}=J_z''\). For any child \(i\) of \(z\),
\(H_{h(i)}\) is computable from \(J_z\), uniformly in \(z,i\): apply
\(r(\mathsf s(a_{z,i}),h(i))\) using column \(i+1\).
Consequently every child oracle is uniformly computable from \(H_{h(z)}\).
\end{hblemma}

\begin{hbexample}{def:p04-doh-leaf-bound}{A leaf bound is deliberately not minimal}
At a node with no children, every dummy is \(1\), so the cumulative
sequence is \(1,\mathsf s(1),\mathsf s(\mathsf s(1)),\ldots\).
Its limit denotes \(\omega\), and \(h(z)\) denotes \(\omega+2\).
The construction thus gives a uniform valid oracle bound even at a leaf;
it does not claim the least possible ordinal or computational complexity
for a particular recursion. Correctness and effective access, not rank
minimality, are the properties used in the characterization proof.
\end{hbexample}

\begin{hbdefinition}{def:p04-doh-alternation}{An alternating product-tree equation}
Let \(T,U\subseteq\N^{<\N}\) be recursive prefix-closed trees containing
the empty sequence. At a pair \(z=(s,t)\) of equal-length nodes put
\[
 \vartheta(s,t)=1\quad\Longleftrightarrow\quad
 \forall m\bigl[s^\frown m\in T\ \Longrightarrow\
       \exists k\,(t^\frown k\in U\ \land\
                \vartheta(s^\frown m,t^\frown k)=1)\bigr].
\]
The equation assigns a bit by well-founded recursion when its synchronized
product tree is well founded. Its empty-existential case is false, and its
empty-universal case is true. For a recursive family \(T_n,U_n\), use
the numerical code
\[
 \operatorname{node}(n,s,t)=
  1+\langle n,\langle\operatorname{code}(s),\operatorname{code}(t)\rangle\rangle
\]
for a valid equal-length pair, and use \(*=0\) as a common node whose
children are their root pairs. Numbers whose decoded sequences are invalid
or outside the indicated trees have no children. All child and
tree-membership tests are decidable.
\end{hbdefinition}

\begin{hblemma}{thm:p04-doh-alternation}{proved=annex:p04.doh.alternation}{Two-jump evaluation of the alternating equation}
If the common product forest is well founded, the equation has a unique
bit solution. A single effectively compiled oracle program computes
\(\vartheta(s,t)\) from \(H_{h(n,s,t)}\), uniformly in its nodes.
All root bits \(\vartheta(\varnothing,\varnothing)\) are then computed
by one total program from the fixed valid oracle \(H_{h(*)}\).
If \(T\) is well founded and \(U\) has a branch, its root bit is \(1\).
If \(T\) has a branch and \(U\) is well founded, its root bit is \(0\).
\end{hblemma}

\subsubsection{The two inclusions}

\begin{hblemma}{thm:p04-doh-tree-normal-form}{proved=annex:p04.doh.tree-normal-form}{Effective tree normal form}
The function-quantifier convention of
Definition~\ref{def:p04-lightface-analytical-hierarchy} and the
set-quantifier convention used here have effective arithmetical translations
in both directions, preserving \(\Sigma^1_1\) and \(\Pi^1_1\).
Function graphs require arithmetic totality and functionality guards;
set characteristic functions require an arithmetic binary-value guard.
The guard is conjoined under an existential quantifier and made an
antecedent under a universal quantifier.
From a lightface \(\Pi^1_1\) definition of a set \(A\subseteq\N\), one
can effectively construct recursive prefix-closed trees \(T_n\), with
\(n\in A\) exactly when \(T_n\) has no infinite branch. The same
construction converts a lightface existential-set arithmetic definition
into the existence of a branch. No effective selection of a branch is
asserted. This matches the independent tree normal form used in
Theorem~\ref{thm:p11-ktf-tree-normal-form}.
\end{hblemma}

\begin{hbtheorem}{thm:p04-doh-delta-to-hyp}{proved=annex:p04.doh.delta-to-hyp}{Compiling complementary definitions into one jump oracle}
Given a \(\Pi^1_1\) definition of \(A\) and a \(\Pi^1_1\) definition of
\(\N\setminus A\), the construction effectively supplies a natural code
\(a\) and a program index \(d\) such that
\[
 a\in\mathcal O,\qquad \varphi_d^{H_a}=\chi_A.
\]
The promise that the two definitions are complementary is essential:
the compiler need not decide that promise, and correctness or validity is
not asserted for an arbitrary incompatible pair of definitions.
Thus every lightface \(\Delta^1_1\) subset of \(\N\) is hyperarithmetical.
\end{hbtheorem}

\begin{hbtheorem}{thm:p04-doh-hyp-to-delta}{proved=annex:p04.doh.hyp-to-delta}{Arithmetic specification of a unique jump hierarchy}
For each fixed \(a\in\mathcal O\), both \(H_a\) and its complement have
lightface \(\Sigma^1_1\) definitions. This holds uniformly by a finite
formula with the notation code as a numerical parameter, with no claim of
validity on invalid codes. The same is true of every set Turing reducible
to \(H_a\). Hence \(\mathsf{HYP}\subseteq\Delta^1_1\).
\end{hbtheorem}

\begin{hbcorollary}{thm:p04-doh-characterization}{proved=annex:p04.doh.characterization}{The exact hyperarithmetical characterization}
With the actual notation and jump conventions above,
\[
 \{A\subseteq\N:\exists a\in\mathcal O\ (A\leq_T H_a)\}
       =\Delta^1_1.
\]
Numerical parameters are harmless: each fixed finite tuple can be inserted
into a lightface definition or a program. Arbitrary set parameters are not
implicitly allowed. Relative analytical pointclasses and relative ordinal
notation systems require their own explicitly fixed oracle conventions;
they are not the unrelativized \(\mathcal O\) with an arbitrary set
smuggled into its limit sequences.
\end{hbcorollary}

The compiler proof is independent of the finite-level arithmetical hierarchy
and of arithmetic-truth closure ordinals. In particular, a total monadic
extension already shown to be lightface \(\Delta^1_1\) can now be identified
as hyperarithmetical without making that classification an assumption of
the truth construction.


\begin{hbdefinition}{def:p04-wf-if-codes}{Well-founded and ill-founded linear-order codes}
Fix the source's effective coding of computable binary relations and let
\(\mathsf{LO}\) be the codes whose decoded relation is a linear order on its
field.  Put
\[
 \begin{aligned}
  \mathsf{WF}&=\{e\in\mathsf{LO}:R_e
       \text{ well-orders its field}\},\\
  \mathsf{IF}&=\{e\in\mathsf{LO}:R_e
       \text{ has an infinite descending sequence}\}.
 \end{aligned}
\]
Thus \(\mathsf{WF}\) and \(\mathsf{IF}\) are complementary within
\(\mathsf{LO}\).  Invalid or nonlinear relation codes belong to neither set;
this convention prevents a malformed code from being silently called
ill-founded.
\index{well-founded relation code}
\index{ill-founded relation code}
\end{hbdefinition}

\begin{hbtheorem}{thm:p04-wf-if-completeness}{cited=Moschovakis2016Hyperarithmetical}{Complete analytical examples}
Under the fixed effective coding, \(\mathsf{WF}\) is \(\Pi^1_1\)-complete and
\(\mathsf{IF}\) is \(\Sigma^1_1\)-complete under effective many-one
reductions whose outputs belong to \(\mathsf{LO}\).
\end{hbtheorem}

\begin{table}[htbp]
\centering
\small
\begin{tabular}{@{}p{.15\textwidth}p{.27\textwidth}p{.22\textwidth}p{.27\textwidth}@{}}
\toprule
Level & Allowed normal form & Calibrating example & Exact evidence \\
\midrule
arithmetical & number quantifiers only & every arithmetical set is in
\(\Delta^1_1\) & \cite[secs.~3A--3D]{Moschovakis2016Hyperarithmetical} \\
\(\Sigma^1_1\) & \(\exists\alpha\,\varphi(n,\alpha)\), arithmetical
\(\varphi\) & \(\mathsf{IF}\) & \(\Sigma^1_1\)-complete;
Theorem~\ref{thm:p04-wf-if-completeness} \\
\(\Pi^1_1\) & \(\forall\alpha\,\varphi(n,\alpha)\), arithmetical
\(\varphi\) & \(\mathsf{WF}\) & \(\Pi^1_1\)-complete;
Theorem~\ref{thm:p04-wf-if-completeness} \\
\(\Delta^1_1\) & both level-one forms & \(\mathsf{HYP}\) & equality by
Theorem~\ref{thm:p04-hyp-delta-one-one} \\
\(\Sigma^1_n,\Pi^1_n\), \(n>1\) & alternating function-quantifier blocks &
\(\mathsf U^1_{\Sigma_n}\), \(\mathsf U^1_{\Pi_n}\) & complete under
effective many-one reductions; Theorem~\ref{thm:p04-universal-analytical-completeness} \\
\bottomrule
\end{tabular}
\caption{The lightface analytical hierarchy.  Function
parameters would change the pointclass and are not tacitly allowed.}
\label{tab:p04-analytical-hierarchy}
\end{table}

\begin{figure}[htbp]
\centering
\begin{tikzpicture}[
  node distance=11mm and 16mm,
  box/.style={draw,rounded corners,minimum width=27mm,minimum height=8mm,
              align=center,font=\small},
  inc/.style={-{Latex[length=2mm]},thick}
]
 \node[box] (arith) {arithmetical sets};
 \node[box,above=of arith] (d1)
   {\(\Delta^1_1=\Sigma^1_1\cap\Pi^1_1=\mathsf{HYP}\)};
 \node[box,above left=of d1] (s1) {\(\Sigma^1_1\)\\\(\mathsf{IF}\) complete};
 \node[box,above right=of d1] (p1) {\(\Pi^1_1\)\\\(\mathsf{WF}\) complete};
 \node[box,above=25mm of d1] (d2) {\(\Delta^1_2=\Sigma^1_2\cap\Pi^1_2\)};
 \node[box,above left=of d2] (s2) {\(\Sigma^1_2\)};
 \node[box,above right=of d2] (p2) {\(\Pi^1_2\)};
 \draw[inc] (arith) -- (d1);
 \draw[inc] (d1) -- (s1);
 \draw[inc] (d1) -- (p1);
 \draw[inc] (s1) -- (d2);
 \draw[inc] (p1) -- (d2);
 \draw[inc] (d2) -- (s2);
 \draw[inc] (d2) -- (p2);
 \node[font=\small,above=4mm of s2] {\(\vdots\)};
 \node[font=\small,above=4mm of p2] {\(\vdots\)};
\end{tikzpicture}
\caption{Lightface analytical-hierarchy inclusions.
Arrows assert inclusion only: they do not assert cover relations or encode
uncited strictness.  The \(\mathsf{WF}/\mathsf{IF}\) labels use
Definition~\ref{def:p04-wf-if-codes}; the central equality is
Theorem~\ref{thm:p04-hyp-delta-one-one}.}
\label{fig:p04-analytical-hierarchy}
\end{figure}

\begin{hbwarning}{def:p04-lightface-boldface}{Lightface is not boldface}
The equality \(\Delta^1_1=\mathsf{HYP}\) concerns effective definitions of
subsets of \(\N\) without arbitrary set parameters.  It is not a statement
about the boldface projective hierarchy, and Kleene's \(\mathcal O\) is not a
canonical enumeration of all countable ordinals.
\end{hbwarning}

\section{Algorithmic information}
\label{sec:p04-kolmogorov-complexity}

\begin{hbdefinition}{def:p04-plain-prefix-complexity}{Plain and prefix-free Kolmogorov complexity}
Let \(\{0,1\}^*\) be the finite binary strings, represented numerically by
the computable length-lexicographic bijection with \(\N\).  Thus ``partial
computable'' below is Definition~\ref{def:partial-computable} transported
across that fixed bijection.  A set \(D\subseteq\{0,1\}^*\) is
\emph{prefix-free} when, for distinct \(p,q\in D\), neither string is a prefix
of the other.  A partial machine is prefix-free when its domain is.  Kraft's
inequality then gives
\[
 \sum_{p\in D}2^{-\operatorname{len}(p)}\leq1.
\]
Prefix-freeness is a combinatorial property of a string set, not a condition on
every decompressor; it is needed for self-delimiting descriptions and the
halting probability below \cite{LiVitanyi2008}.

For a partial computable
decompressor \(V:\{0,1\}^*\rightharpoonup\{0,1\}^*\), its \emph{plain
complexity} is
\[
 C_V(x)=\min\{\operatorname{len}(p):V(p)=x\},
\]
where \(\operatorname{len}(p)\) is the number of bits in the finite string
\(p\).
If the domain of \(V\) is prefix-free, its \emph{prefix complexity} is
\[
 K_V(x)=\min\{\operatorname{len}(p):V(p)=x\}.
\]
The minimum is \(\infty\) when there is no description.  The two symbols name
different complexity notions; in particular, prefix-freeness is not a tacit
condition on \(C_V\) \cite{Kolmogorov1965,LiVitanyi2008}.
\index{Kolmogorov complexity!plain}
\index{Kolmogorov complexity!prefix-free}
\index{prefix-free code}
\end{hbdefinition}

\begin{hbtheorem}{thm:p04-kolmogorov-invariance}{cited=LiVitanyi2008}{Invariance theorem}
There is an optimal plain machine \(U_C\) such that, for every partial
computable decompressor \(V\), some constant \(c_V\) satisfies
\[
 C_{U_C}(x)\leq C_V(x)+c_V
 \qquad\text{for every }x\in\{0,1\}^*.
\]
There is likewise a universal optimal prefix-free machine \(U_K\).  More
precisely, for every prefix-free partial computable machine \(V\), there is a
finite binary string \(\sigma_V\) such that
\[
 U_K(\sigma_Vp)\simeq V(p)\qquad\text{for every }p\in\{0,1\}^*,
\]
where juxtaposition denotes concatenation.  The strings \(\sigma_V\) are
chosen from a prefix-free effective coding of machines, so that
\(\operatorname{dom}(U_K)\) remains prefix-free.  Consequently
some constant \(d_V\) satisfies
\[
 K_{U_K}(x)\leq K_V(x)+d_V
 \qquad\text{for every }x\in\{0,1\}^*.
\]
Each constant depends on the two machines but not on \(x\).  Thus complexity
is invariant only up to an additive machine-dependent constant
\cite{Kolmogorov1965}.
\end{hbtheorem}

\begin{hbtheorem}{thm:p04-incompressibility-count}{cited=LiVitanyi2008}{Incompressibility count}
Fix the optimal plain machine \(U_C\).  If \(0\leq c\leq n\), fewer than
\(2^{n-c}\) strings \(x\) of length \(n\) satisfy
\(C_{U_C}(x)<n-c\).  Consequently the fraction of length-\(n\) strings
compressible by at least \(c\) bits is less than \(2^{-c}\).
\end{hbtheorem}

\begin{hbdefinition}{def:p04-upper-semicomputable}{Upper semicomputable function}
An extended-natural-valued function
\(h:\{0,1\}^*\to\N\cup\{\infty\}\) is \emph{upper semicomputable} when
\[
 \{(x,n):h(x)<n\}
\]
is computably enumerable.  Equivalently, there is a total computable sequence
\(h_s(x)\in\N\cup\{\infty\}\), using one fixed code for \(\infty\), that is
nonincreasing in \(s\) and converges to \(h(x)\).  This is an approximation
property of a numerical function, not enumerability of its graph or of a set;
it is strictly weaker than computability and is used in the next theorem and
the information table.
\index{upper semicomputable function}
\end{hbdefinition}

\begin{hbtheorem}{thm:p04-kolmogorov-noncomputable}{cited=LiVitanyi2008}{Computability boundary for complexity}
For optimal effective machines, plain complexity and prefix complexity are
upper semicomputable but not computable.
\end{hbtheorem}

\begin{hbtheorem}{thm:p04-chaitin-incompleteness}{cited=Chaitin1974}{Chaitin-style incompleteness}
Fix the optimal prefix-free machine \(U_K\), an arithmetic coding of strings and
programs.  Let \(\mathsf{Run}_{U_K}(p,x,t)\) be the bounded arithmetic
formula asserting that \(t\) codes a valid halting computation of \(U_K\) on
\(p\) with output \(x\).  This bounded arithmetic formula is needed because
\(U_K\) uses the separately fixed prefix-machine and string coding.  It is an
object-language computation predicate, not a machine relation; unlike the
inherited \(\mathsf{Comp}(e,x,s)\), its machine index is fixed and its input
is a binary program string.  A primitive-recursive compiler between the two
finite trace codings preserves acceptance, and
\(\mathsf{Run}_{U_K}\) is used only in the following theory hypothesis and
complexity sentence.

Let \(T\) be a computably axiomatized arithmetic
theory with a primitive-recursive proof relation such that every true closed
instance of \(\mathsf{Run}_{U_K}\) is provable in \(T\), and interpret
\(K_{U_K}(x)>n\) by the \(\Pi^0_1\) sentence saying that no program of length
at most \(n\) has such a computation.  Assume \(T\) is \(\Pi^0_1\)-sound:
every \(\Pi^0_1\) sentence it proves is true in \(\N\).

Then there is a constant \(b_{T,U_K}\) such that \(T\) proves no sentence
\(K_{U_K}(x)>n\) with \(n>b_{T,U_K}\).  The bound depends on \(T\), \(U_K\),
the proof relation, and the coding; mere consistency is not substituted for
the displayed soundness and computation-representation hypotheses.
\end{hbtheorem}

\begin{hbdefinition}{def:p04-chaitin-omega}{Prefix-free halting probability}
For the fixed universal prefix-free machine \(U_K\), put
\[
 \Omega_{U_K}=\sum_{U_K(p)\mathord\downarrow}2^{-\operatorname{len}(p)}.
\]
Prefix-freeness and Kraft's inequality make this convergent and identify it
with the fair-coin measure of the halting domain.  A real \(r\) is
\emph{left computably enumerable} (\emph{left-c.e.}) when a computable nondecreasing sequence of
rationals has supremum \(r\).  This is a one-sided approximation property of a
real, not enumerability of its binary digits.  It is needed to distinguish the
effective lower approximation to \(\Omega_{U_K}\) from computability of the
real itself.  The definition depends on the fixed prefix-free universal
machine and does not attach \(\Omega\) to an arbitrary plain machine
\cite{Chaitin1975,LiVitanyi2008}.
\index{Chaitin halting probability}
\end{hbdefinition}

\begin{hbtheorem}{thm:p04-omega-halting}{proved=annex:p4.frontiers.2}{Information carried by \(\Omega\)}
The real \(\Omega_{U_K}\) is left computably enumerable and noncomputable.
For each \(n\), its first \(n\) binary digits determine which programs of
length at most \(n\) halt on \(U_K\).
\end{hbtheorem}

The annex proves the result from the increasing halting-weight approximation
and the fixed machine's undecidable halting domain.  No broader randomness
calculus is asserted here \cite{Chaitin1975,LiVitanyi2008}.

\begin{table}[htbp]
\centering
\small
\begin{tabular}{@{}p{.17\textwidth}p{.24\textwidth}p{.25\textwidth}p{.25\textwidth}@{}}
\toprule
Object & Machine requirement & Licensed conclusion & Boundary \\
\midrule
\(C_{U_C}(x)\) & optimal plain decompressor & invariant up to an additive
constant; upper semicomputable, not computable & no prefix-free probability
claim \\
\(K_{U_K}(x)\) & optimal prefix-free machine & same invariance form; Chaitin
bound with the stated theory-soundness hypothesis & values still depend on
the chosen machine \\
\(\Omega_{U_K}\) & universal prefix-free machine & left-c.e. noncomputable
halting measure; first \(n\) bits decide halting through length \(n\) & no
general randomness chapter \\
\bottomrule
\end{tabular}
\caption{Plain and prefix-free algorithmic information.
The table prevents the invariance theorem from being read as literal machine
independence.}
\label{tab:p04-kolmogorov-boundaries}
\end{table}

\section{Hilbert's tenth problem}
\label{sec:p04-mrdp}

\begin{hbdefinition}{def:p04-diophantine-set}{Diophantine set}
A set \(A\subseteq\N^k\) is \emph{Diophantine} when there are
\(m\in\N\) and a polynomial \(p\) with integer coefficients such that
\[
 \vec x\in A\quad\Leftrightarrow\quad
 \exists y_1,\ldots,y_m\in\N\;
 p(\vec x,y_1,\ldots,y_m)=0.
\]
This is a class of sets represented by existential polynomial equations, not
a complexity class.  When an instance is instead posed over \(\mathbb Z\),
the domain is stated explicitly.  Integer variables can be represented as
differences of naturals, and natural variables can be represented over
\(\mathbb Z\) by sums of four squares; these are domain translations, not a
license to suppress the domain \cite{Davis1973}.
\index{Diophantine set}
\end{hbdefinition}

\begin{hbtheorem}{thm:p04-mrdp}{cited=Davis1973}{Diophantine characterization of c.e. sets}
For every \(k\geq1\), a set \(A\subseteq\N^k\) is Diophantine if and only if
it is computably enumerable.  The difficult implication is the
Matiyasevich--Robinson--Davis--Putnam (MRDP) theorem
\cite{DavisPutnamRobinson1961,Matiyasevich1970}.
\end{hbtheorem}

\begin{hbcorollary}{thm:p04-hilbert-tenth}{cited=Davis1973}{Undecidability of Diophantine solvability}
There is no algorithm that, given an integer-coefficient polynomial in
finitely many variables, decides whether it has a solution in integers.
Equivalently, after the explicit domain translations in
Definition~\ref{def:p04-diophantine-set}, there is no such algorithm for
solutions in natural numbers.
\end{hbcorollary}

\begin{hbmethod}{def:p04-mrdp-proof-survey}{MRDP proof survey}
The easy direction enumerates tuples and checks polynomial equality.
Davis--Putnam--Robinson first obtained exponential-Diophantine
representations of c.e. sets.  Julia Robinson isolated a growth condition
sufficient to make exponentiation Diophantine, and Matiyasevich supplied that
condition using rapidly growing recurrence sequences.

The Pell core fixes \(a>1\), puts \(d=a^2-1\), and studies
\[
 x^2-dy^2=1,
 \qquad
 x_n(a)+y_n(a)\sqrt d=(a+\sqrt d)^n.
\]
Every nonnegative integer solution is one pair
\((x_n(a),y_n(a))\), and multiplication gives the recurrences
\[
\begin{aligned}
 x_{m+n}&=x_mx_n+d\,y_my_n,\\
 y_{m+n}&=x_my_n+x_ny_m.
\end{aligned}
\]
The cited Pell-sequence lemmas give Diophantine definitions of the required
solution and congruence relations.  Julia Robinson's growth condition asks for
a Diophantine set \(D\) of pairs \((u,v)\) such that every pair satisfies
\(v<u^u\), while for every \(k\) some pair satisfies \(v>u^k\).  Her theorem
turns such a set into a Diophantine definition of exponentiation.
Matiyasevich proved the condition using a Fibonacci-type recurrence whose
even-index subsequence is Diophantine.  Once exponentiation is removed from
the Davis--Putnam--Robinson exponential-Diophantine representation, every
c.e. set has an ordinary polynomial Diophantine representation
\cite{DavisPutnamRobinson1961,Matiyasevich1970,Davis1973}.

This is a proof survey, not the full MRDP proof.  It omits the explicit
polynomials coding the Pell solution index and congruences, the Chinese-
remainder elimination of bounded universal quantifiers, and the final
composition of those polynomial representations; those are the substantive
lemmas supplied by the cited proofs.
\index{Hilbert's tenth problem}
\end{hbmethod}

\section{Undecidability catalog}
\label{sec:p04-undecidability-catalog}

\begin{hbconvention}{def:p04-finite-instance-codings}{Effective encodings for the catalog}
Encode an integer polynomial by its finite list of pairs consisting of a signed
binary coefficient and an exponent vector; encode a finite semigroup or group
presentation by a finite generator alphabet followed by a finite list of word
pairs; and encode a Wang-tile set by a finite list of four-tuples of natural-
number color codes in north--east--south--west order.  Use self-delimiting
length prefixes for every finite list and word.  Syntactic validity, decoding,
and re-encoding are computable.  Any two standard effective encodings of these
finite objects are computably intertranslatable, so the undecidability results
do not depend on this particular serialization.  This convention fixes input
codes only; each row below separately fixes its yes predicate.
\index{encoding!undecidable problems}
\end{hbconvention}

\begin{hbdefinition}{def:p04-fixed-word-problem}{Word problem for a fixed presentation}
For one fixed finite semigroup presentation \(\langle X\mid R\rangle\), the
instances are pairs of words over \(X\), and the yes predicate is equality of
the elements represented by those words.  For a fixed finite group
presentation the instances and yes predicate are analogous, with group-word
syntax.  These fixed-presentation problems are stronger existence results
than the uniform problem whose input also supplies a presentation.
\index{word problem!semigroup}
\index{word problem!group}
\end{hbdefinition}

\begin{hbtheorem}{thm:p04-semigroup-word-problem}{cited=Post1947}{Post--Markov semigroup word-problem theorem}
There exists a finitely presented semigroup whose fixed-presentation word
problem is undecidable \cite{Post1947,Markov1947}.
\end{hbtheorem}

\begin{hbtheorem}{thm:p04-group-word-problem}{cited=Boone1959}{Novikov--Boone theorem}
There exists a finitely presented group whose fixed-presentation word problem
is undecidable \cite{Novikov1955,Boone1958,Boone1959}.
\end{hbtheorem}

\begin{hbdefinition}{def:p04-domino-problem}{Whole-plane domino problem}
An instance is a finite set of unit-square Wang tile types with colored
edges.  Unlimited translated copies may be used, but rotation is forbidden.
The yes predicate holds exactly when the entire integer plane can be tiled so
that every pair of adjacent edges has matching colors.  This is not a problem
about tiling a supplied finite board \cite{Wang1961PatternII}.
\index{domino problem}
\index{Wang tile}
\end{hbdefinition}

\begin{hbtheorem}{thm:p04-domino-undecidable}{cited=Berger1966}{Undecidability of the domino problem}
The whole-plane domino problem in
Definition~\ref{def:p04-domino-problem} is undecidable.
\end{hbtheorem}

\begingroup\small
\begin{longtable}{@{}p{0.19\textwidth}p{0.40\textwidth}p{0.32\textwidth}@{}}
\caption{Exact undecidability catalog.  Every row fixes an
encoded instance and a yes predicate.}
\label{tab:p04-undecidability-catalog}\\
\toprule
Problem & Encoded instance and yes predicate & Result and evidence \\
\midrule
\endfirsthead
\toprule
Problem & Encoded instance and yes predicate & Result and evidence \\
\midrule
\endhead
\bottomrule
\endfoot
Self-application domain & A program index \(e\); yes exactly when
\(\varphi_e(e)\mathord\downarrow\) under
the indexed-program definedness convention in
Section~\ref{sec:universal-reductions-01}. & c.e. and undecidable;
Theorem~\ref{thm:halting-undecidable}; \cite{Turing1937}. \\
\addlinespace
Rice index problem & A program index \(e\); yes exactly when the unary
partial function computed by \(e\) has one fixed nontrivial extensional
property. & undecidable; Theorem~\ref{thm:rice}; \cite{Rice1953}.
Syntactic properties are outside the theorem. \\
\addlinespace
Hilbert's tenth problem & An integer-coefficient polynomial; yes exactly when
it has an integer solution. & undecidable;
Corollary~\ref{thm:p04-hilbert-tenth};
\cite{DavisPutnamRobinson1961,Matiyasevich1970,Davis1973}. \\
\addlinespace
Semigroup word problem & Two words for one fixed finite semigroup
presentation; yes exactly when they represent the same element. & undecidable
for some fixed presentation; Theorem~\ref{thm:p04-semigroup-word-problem};
\cite{Post1947,Markov1947}. \\
\addlinespace
Group word problem & Two words for one fixed finite group presentation; yes
exactly when they represent the same element. & undecidable for some fixed
presentation; Theorem~\ref{thm:p04-group-word-problem};
\cite{Novikov1955,Boone1958,Boone1959}. \\
\addlinespace
Domino/Wang tiling & A finite Wang-tile set; yes exactly when translated,
unrotated copies tile the whole integer plane with matching adjacent colors. &
undecidable; Theorem~\ref{thm:p04-domino-undecidable};
\cite{Berger1966,Wang1961PatternII}. \\
\addlinespace
First-order validity & A code of a first-order sentence in the fixed effective
syntax; yes exactly when it is valid in every structure of the stated
signature. & c.e. and undecidable;
Theorem~\ref{thm:fol-validity-undecidable} and
Corollary~\ref{thm:fol-validity-enumerable};
\cite{Church1936,Turing1937}.  Restricted
signatures require separate results. \\
\end{longtable}
\endgroup

\section{Busy Beaver and computable presentations}
\label{sec:p04-busy-beaver-structures}

\begin{hbdefinition}{def:p04-busy-beaver}{Busy Beaver functions under a fixed machine convention}
For this definition only, an \(n\)-state machine has \(n\geq1\) nonhalting
states, one separate halting state, alphabet \(\{0,1\}\) with blank \(0\), a
two-way infinite tape, and one left or right move per transition.  It starts
in a designated nonhalting state on an all-zero tape.  The transition into
the halting state counts as one step, and the halting state has no outgoing
transition.

Among the finitely many such machines that halt, let
\(\operatorname{BB}_{\mathrm{step}}(n)\) be the largest number of executed
transitions and let \(\operatorname{BB}_{1}(n)\) be the largest number of
tape cells containing \(1\) at halt.  These functions belong to this exact
state, symbol, blank, motion, and step convention; their small numerical
values are not transferred to another convention without a translation.
\index{Busy Beaver function}
\end{hbdefinition}

\begin{hbtheorem}{thm:p04-busy-beaver-noncomputable}{cited=Rado1962BusyBeaver}{Rado's noncomputability theorem}
The Busy Beaver functions \(\operatorname{BB}_{\mathrm{step}}\) and
\(\operatorname{BB}_{1}\), under the fixed finite-machine normalization
above, are not computable.
\end{hbtheorem}

The step-bound conclusion is historically tied to the blank-tape halting
problem, while the productivity conclusion for
\(\operatorname{BB}_{1}\) is a separate clause of Rado's theorem.  Both are
imported under the displayed citation status \cite{Rado1962BusyBeaver}.

The definition of a computable presentation and the distinction between one
computable copy and computable categoricity are
Definitions~\ref{def:p03-computable-presentation}
and~\ref{def:p03-computably-categorical}.  Fix the handbook's program
numbering and a computable serialization of atomic sentences.  An
\emph{index of a computable presentation} is an index \(e\) whose program is
total and decides the serialized atomic diagram of that presentation.  The
index is coding-dependent metalinguistic data, not an isomorphism invariant of
the abstract structure; it differs from the presentation itself and is used by
the uniform corollary below.
\index{computable presentation!index}

\begin{hbcorollary}{thm:p04-rational-order-uniform-index}{proved=annex:p4.frontiers.1}{Uniform effective back-and-forth}
There is a total computable compiler \(F(e,d)\) such that, whenever \(e,d\)
index computable presentations of the countable dense linear order without
endpoints, \(F(e,d)\) indexes a computable isomorphism between those
presentations.
\end{hbcorollary}

Thus the complete Part III categoricity result has a uniform index-level
consequence; no claim is made that every countable structure has a computable
copy.  See \textcite[chaps.~1--3]{AshKnight2000} for the broader theory.

\begin{hbwarning}{def:p04-frontier-ceilings}{Frontier claim ceilings}
The pointclass results above are lightface; the MRDP discussion is a proof
survey rather than a full proof; the word problems use fixed finite
presentations; the domino problem concerns the whole plane; Busy Beaver values
depend on the machine convention; and the computable-structure entry is a
uniform corollary over valid presentation indices.  None of these qualifications may be dropped when the results are
reused in a later table.
\end{hbwarning}

%% Coverage: D-P5-006, D-P9-002, N-a-019--024, A-046, A-048, A-125.
